% !TeX root = ../navier-stokes-manuscript.tex
% ============================================================
% 1.7 Singularity and the Derivative Tower
% ============================================================
\subsection{Singularity and the Derivative Tower}\label{sub:SingularityAndTheDerivativeTower}

A finite limiting velocity along one material parcel does not settle whether the velocity field extends smoothly to that time.
Let \(X(a,t)\) be the trajectory of the parcel labelled by \(a\), and write its carried velocity as
\[
  \dot X(a,t)=u(t,X(a,t)),
  \qquad
  U(a,t):=u(t,X(a,t)).
\]
The curve \(U(a,\cdot)\) is the velocity that the parcel actually experiences.

A simple time profile shows how much its terminal value can hide.
Fix a unit vector \(e\) and suppose, kinematically, that near a finite time \(T\),
\[
  U(t)=U_*+(T-t)^{3/2}e,
  \qquad t<T.
\]
Then
\[
  U(t)\longrightarrow U_*,
  \qquad
  \dot U(t)=-\frac32(T-t)^{1/2}e\longrightarrow0,
  \qquad
  |\ddot U(t)|=\frac34(T-t)^{-1/2}|e|\longrightarrow\infty.
\]
The velocity arrives, and even the acceleration tends to zero, while the jerk diverges.
This is a kinematic profile, with no claim that a \NavierStokes{} solution realizes it.
It isolates one possible kinematic regularity failure: the carried velocity and its acceleration can approach finite limits while the jerk does not.

The \NavierStokes{} equation already contains the parcel's first response.
Set
\[
  D_t:=\partial_t+u\cdot\nabla.
\]
The chain rule along \(X(a,t)\) and the momentum equation give
\[
  \dot U(a,t)
  =
  D_tu(t,X(a,t))
  =
  \bigl(-\nabla p+\nu\Delta u\bigr)(t,X(a,t)).
\]
The pressure force and the viscous response to spatial curvature determine how that same parcel accelerates.
Differentiate the same equality along the same path, and its jerk is
\[
  \ddot U(a,t)
  =
  D_t^2u(t,X(a,t))
  =
  D_t\bigl(-\nabla p+\nu\Delta u\bigr)(t,X(a,t)).
\]
Another material derivative gives the rate at which the jerk changes, and the process continues at every finite order while the field remains smooth.
Higher names the derivative order.
The response may shrink, pass through zero, or vanish.
Each material rung is the derivative of the one below, evaluated from the same \NavierStokes{} field.

Material differentiation follows the parcel, while Eulerian differentiation watches the field at one fixed spatial coordinate.
Write the fixed-coordinate derivatives as
\[
  V_n:=\partial_t^n u.
\]
Expanding the material jerk gives
\[
  D_t^2u
  =
  V_2
  +2(V_0\cdot\nabla)V_1
  +(V_1\cdot\nabla)V_0
  +(V_0\cdot\nabla)\bigl((V_0\cdot\nabla)V_0\bigr).
\]
The fixed-position second time derivative \(V_2\) is one part of the jerk carried by the parcel.
The remaining terms record the parcel moving through spatial variation while the field changes in time.
The identity \(D_t=\partial_t+u\cdot\nabla\) keeps the two readings joined at every order.

\divider{The Temporal Rungs}

The parcel has led us from velocity to material acceleration and jerk.
To expose the fixed-coordinate derivatives from which those carried responses are assembled, differentiate the same field equation in time while holding \(x\) fixed.
Write the pressure derivatives as
\[
  Q_n:=\partial_t^n p.
\]
The original momentum equation gives the first Eulerian rung:
\[
  V_1
  =
  \nu\Delta V_0
  -(V_0\cdot\nabla)V_0
  -\nabla Q_0.
\]
The time change of the velocity remains tied to transport, pressure, and viscous curvature.
Differentiate once in time.
The derivative can land on either velocity factor in the transport term, which gives
\[
  V_2
  =
  \nu\Delta V_1
  -(V_1\cdot\nabla)V_0
  -(V_0\cdot\nabla)V_1
  -\nabla Q_1.
\]
The new row contains both ways the transporting field and the transported field can change.
Differentiate again.
The middle allocation now occurs twice:
\[
  V_3
  =
  \nu\Delta V_2
  -(V_2\cdot\nabla)V_0
  -2(V_1\cdot\nabla)V_1
  -(V_0\cdot\nabla)V_2
  -\nabla Q_2.
\]
One more time derivative gives three copies of each neighbouring split:
\[
  \begin{aligned}
  V_4
  ={}&
  \nu\Delta V_3
  -(V_3\cdot\nabla)V_0
  -3(V_2\cdot\nabla)V_1\\
  &-3(V_1\cdot\nabla)V_2
  -(V_0\cdot\nabla)V_3
  -\nabla Q_3.
  \end{aligned}
\]
The coefficients \(1,1\), then \(1,2,1\), then \(1,3,3,1\), count every way the time derivatives can split between the two copies of velocity.
The first four rungs have exposed the same Pascal pattern.
At arbitrary order it reads
\[
  V_{n+1}
  +
  \sum_{k=0}^{n}
  \binom{n}{k}
  (V_k\cdot\nabla)V_{n-k}
  +
  \nabla Q_n
  =
  \nu\Delta V_n,
  \qquad
  \nabla\cdot V_n=0.
\]
Every temporal rung is a fixed-coordinate derivative of the same velocity and pressure field.
The recurrence couples it to lower rows through transport, to spatial curvature through viscosity, and to pressure at the same time order.

\divider{Pressure and the Spatial Rungs}

The temporal recurrence already contains \(\nabla Q_n\), \(\nabla V_n\), and \(\Delta V_n\).
On the classical interval, fixed-coordinate time differentiation commutes with the spatial operators, so the pressure gradient and spatial derivatives remain in every temporal row.
Taking the divergence of the original momentum equation and using \(\nabla\cdot u=0\) gives
\[
  -\Delta p
  =
  \partial_i\partial_j(u_i u_j),
\]
where repeated spatial indices are summed.
After \(n\) time derivatives,
\[
  -\Delta Q_n
  =
  \partial_i\partial_j
  \sum_{k=0}^{n}
  \binom{n}{k}
  V_{k,i}V_{n-k,j}.
\]
The elliptic equation determines \(Q_n\) up to a spatial constant; the manuscript's decay and pressure normalization fix that constant.
The pressure response at one location can depend on the velocity configuration far beyond an arbitrarily small neighbourhood of that location.

Spatial differentiation opens the same momentum law in another direction.
Differentiate once in the coordinate \(x_\ell\):
\[
  \begin{aligned}
  \partial_t(\partial_\ell u)
  &+
  (u\cdot\nabla)(\partial_\ell u)
  +
  ((\partial_\ell u)\cdot\nabla)u
  +
  \nabla(\partial_\ell p)\\
  &=
  \nu\Delta(\partial_\ell u).
  \end{aligned}
\]
The new transport term records how the spatial variation of the transporting velocity acts on the original velocity gradient.
At the next spatial rung, the product rule creates both cross terms:
\[
  \begin{aligned}
  \partial_t(\partial_{\ell m}u)
  &+
  (u\cdot\nabla)(\partial_{\ell m}u)
  +
  ((\partial_\ell u)\cdot\nabla)(\partial_m u)\\
  &+
  ((\partial_m u)\cdot\nabla)(\partial_\ell u)
  +
  ((\partial_{\ell m}u)\cdot\nabla)u
  +
  \nabla(\partial_{\ell m}p)\\
  &=
  \nu\Delta(\partial_{\ell m}u).
  \end{aligned}
\]
Every further time or spatial derivative repeats this nesting and creates the corresponding mixed row.

Using a multi-index \(\alpha\) for the spatial derivatives, collect the velocity coordinates at one spacetime point in
\[
  \mathcal J^\infty u(t,x)
  :=
  \bigl(
    \partial_t^n\partial_x^\alpha u(t,x)
  \bigr)_{n\geq0,\,\alpha\in\mathbb N_0^3}.
\]
This is the derivative tower of the velocity.
On every classical part of the solution, the tower is already present; following the parcel and differentiating the law exposes the relations among its rungs.
At each classical instant, its rungs are simultaneous readings of one velocity field.
Material acceleration, jerk, and every higher carried response are combinations of these fixed-coordinate derivatives along the trajectory.
The pressure rows stay attached through the differentiated momentum equations, the Poisson equations, and the divergence-free constraint.

Individual rungs are allowed to vanish.
The zero solution has a vanishing tower, and any smooth steady flow has \(V_n=0\) for every \(n\geq1\); these vanishing temporal rungs are compatible with smoothness.
Neither zero nor nonzero decides lawfulness.
A rung is lawful when its value, together with the other terms at that order, satisfies the corresponding differentiated equation.

A finite terminal value of \(U\) leaves the higher tower unsettled.
Every derivative required for a smooth endpoint must extend to \(T\) with the relevant finite and compatible control.
Failure of a temporal, spatial, or mixed rung is one way smooth extension can break down; this diagnostic does not assert that every loss of continuation has a single pointwise first rung, or that pointwise jet limits alone suffice for continuation.
